\section{The low-defect projection in the original auxiliary sectors}
\label{sec:replica_low_defect_projection}

Let $A,C,R$ be finite coordinate sets, with associated physical and
auxiliary Hilbert spaces $\mathcal H_A,\mathcal H_C,\mathcal H_R$.
For $k\ge0$, all operators below act on
$\mathcal H_A^{\otimes k}\otimes
(\mathcal H_C^{\otimes k}\otimes\mathcal H_R^{\otimes k})$.
The coordinate sets may be empty, and a zeroth tensor power is $\mathbb C$.

\begin{definition}[Simultaneous copy action]
\label{def:replica_joint_copy_action}
\lean{TensorPower.replicaJointCopyPerm}
\leanok
A permutation $\sigma\in S_k$ acts simultaneously on the physical
copy register and the two auxiliary copy registers.
\end{definition}

\begin{theorem}[The simultaneous permutation matrix]
\label{thm:replica_joint_copy_operator}
\lean{TensorPower.permOp_replicaJointCopyPerm}
\leanok
\uses{def:replica_joint_copy_action}
The matrix of this action is
$U^A_\sigma\otimes(U^C_\sigma\otimes U^R_\sigma)$.
\end{theorem}
\begin{proof}
\leanok
Both matrices have entry one precisely when each input copy string is
sent to its corresponding output string, and entry zero otherwise.
\end{proof}

\begin{definition}[The low-defect operator]
\label{def:replica_low_defect_projection}
\lean{Matrix.replicaLowDefectProjection}
\leanok
\uses{def:replica_joint_copy_action}
For a physical vector $\Omega$, let
$N_\Omega=\sum_{i=0}^{k-1}(I-|\Omega\rangle\langle\Omega|)_i$.
Fix the two original auxiliary labels $\ell_C,\ell_R$ and $\tau\in\mathbb R$.
Write $E_C,E_R$ for their whole-copy label projections, and set
\begin{align}
 Q&=\mathbf1_{(-\infty,\tau k]}(N_\Omega)\otimes I_{CR},\\
 L&=I_A\otimes(E_C\otimes E_R),\\
 P&=QLS,
\end{align}
where $S$ is the average of the simultaneous permutation matrices.
Here the identity operators act on the corresponding copy registers.
This is the low-defect construction in
\cite[Section~8, \textup{(comparator:defect-mass)}]{OpenAI2026AreaLaw}.
\end{definition}

\begin{theorem}[Orthogonality and retained conditions]
\label{thm:replica_low_defect_projection_conditions}
\lean{Matrix.isStarProjection_replicaLowDefectProjection,Matrix.replicaLowDefectProjection_fixed_conditions}
\leanok
\uses{def:replica_low_defect_projection}
If $\|\Omega\|=1$, then $P$ is an orthogonal projection.
Every vector $u$ satisfying $Pu=u$ lies in the range of $Q$, has both
original auxiliary labels, and is fixed by every simultaneous copy
permutation. No commutation with a replica metric is assumed.
\end{theorem}
\begin{proof}
\leanok
\uses{thm:replica_cfc_auxiliary_commute,lem:label_projectors,thm:replica_joint_copy_operator}
The physical cutoff commutes with every physical copy permutation and
acts trivially on the auxiliary registers. Each auxiliary label
projection commutes with its copy action. Thus $Q,L,S$ are commuting
orthogonal projections. Their product is an orthogonal projection
whose range lies in each of their ranges. The range of $L$ lies in
each of the two original auxiliary label sectors.
\end{proof}

\begin{theorem}[Action on the original sectors]
\label{thm:replica_low_defect_projection_on_sector}
\lean{Matrix.replicaLowDefectProjection_mulVec_of_sector}
\leanok
\uses{def:replica_low_defect_projection}
For arbitrary $\Omega$, if $u$ has both original auxiliary labels and
simultaneous copy symmetry, then $Pu=Qu$.
\end{theorem}
\begin{proof}
\leanok
The symmetry average and each auxiliary label projection fix $u$.
Apply the factors of $P$ in their defining order.
\end{proof}

\begin{theorem}[Retained mass from the physical gap]
\label{thm:replica_low_defect_projection_gap_mass}
\lean{Matrix.replicaLowDefectProjection_gap_mass_ge}
\leanok
Suppose $\|\Omega\|=1$, $g>0$, $k>0$, $\tau>0$, and
\begin{align}
 H-E_0I\ge g(I-|\Omega\rangle\langle\Omega|).
\end{align}
For every unit vector $v$ with simultaneous copy symmetry and both
original auxiliary labels,
\begin{align}
 \langle v,Pv\rangle\ge
 1-\frac{\langle v,(\overline H\otimes I_{CR})v\rangle-E_0}{g\tau},
 \qquad \overline H=\frac1k\sum_{i=0}^{k-1}H_i.
\end{align}
This is the retained-mass estimate in
\cite[Section~8, \textup{(comparator:defect-mass)}]{OpenAI2026AreaLaw}.
\end{theorem}
\begin{proof}
\leanok
\uses{thm:replica_low_defect_projection_on_sector,thm:replica_cutoff_mass}
On $v$, the projection $P$ acts as the physical cutoff $Q$. The stated
inequality is therefore the spectral first-moment bound from the
physical gap.
\end{proof}

\begin{theorem}[Compression from component bounds]
\label{thm:replica_low_defect_compression_from_components}
\lean{Matrix.replicaLowDefectProjection_compression_le_of_component_bounds}
\leanok
Let $\|\Omega\|=1$, $0\le\tau\le\tfrac12$, $c\ge0$, and $H\ge0$.
For $B\subseteq\{0,\ldots,k-1\}$, let $R_B$ be the physical projection
onto precisely the copies in $B$ being orthogonal to $\Omega$.
Suppose that, for every vector $u$ with simultaneous copy symmetry
and both original auxiliary labels and every $B$ with $|B|\le\tau k$,
the components $w_B=(R_B\otimes I_{CR})u$ satisfy
\begin{align}
 \langle w_B,Hw_B\rangle\le c\|w_B\|^2.
\end{align}
Then
\begin{align}
 PHP\le(k+1)e^{kh(\tau)}cP,
\end{align}
where $h(\tau)=-\tau\log\tau-(1-\tau)\log(1-\tau)$, with
$0\log0=0$. The component estimate remains an explicit hypothesis;
this is the subsequent implication used in
\cite[Section~8, \textup{(comparator:inverse-compression)}]{OpenAI2026AreaLaw}.
\end{theorem}
\begin{proof}
\leanok
\uses{thm:replica_low_defect_projection_conditions,thm:replica_excitation_compression_bound}
Every vector in the range of $P$ has the required cutoff, symmetry,
and label properties. Apply the quadratic-form estimate for its sum
of excitation components and the bound
$(k+1)e^{kh(\tau)}$ for the number of retained subsets.
\end{proof}
